\documentclass[11pt,a4paper]{article}

\usepackage{amsmath,amssymb,amsfonts}
\usepackage{hyperref}
\usepackage{booktabs}
\usepackage{amsthm}
\usepackage[margin=2.5cm]{geometry}

\newtheorem{theorem}{Theorem}
\newtheorem{proposition}[theorem]{Proposition}

\newcommand{\Spec}{\mathrm{Spec}}
\newcommand{\Zp}{\mathbb{Z}[1/p]}
\newcommand{\lP}{\ell_{\mathrm{P}}}

\begin{document}

\title{Can Spacetime Source Code Editing Crash the Universe?\\
  A Safety Analysis of Arithmetic Vacuum Engineering}

\author{Wright Brothers\\{\small Independent Research}}
\date{April 2026}
\maketitle

\begin{abstract}
If spacetime has the arithmetic structure of $\Spec(\mathbb{Z})$
and physical law can be locally modified by localization
$\mathbb{Z} \to \Zp$ (``prime muting''), a natural concern arises:
could such modifications trigger uncontrolled vacuum decay,
consuming the universe? We perform a rigorous safety analysis.
Our central result is a no-go theorem: the Coleman-De Luccia
bubble nucleation mechanism, which drives vacuum decay in standard
QFT, is \emph{topologically forbidden} when the two vacua have
different K-theory invariants ($K_1(\mathbb{Z}) = \mathbb{Z}/2$
vs.\ $K_1(\Zp) = \mathbb{Z}/2 \times \mathbb{Z}$). No continuous
path in field space connects them, so no bounce solution exists and
the nucleation rate is exactly zero. We identify five safety
mechanisms (topological protection, wall radiation, Planck-scale
minimum size, the Casimir precedent, and parity reversibility),
address five counter-arguments (finite temperature, quantum
tunneling, cosmic rays, black holes, loss of control), and
propose a concrete BEC sandbox protocol for pre-deployment
safety testing. We show that the safety of arithmetic vacuum
engineering rests primarily on the topological no-go theorem ---
the wall radiation and other mechanisms provide backup but are
individually insufficient for large bubbles. We conclude that
the concern of universal collapse is theoretically unfounded, but
recommend the BEC sandbox protocol as a precautionary measure.
\end{abstract}

% ============================================================================
\section{Introduction}
% ============================================================================

The arithmetic vacuum engineering program~\cite{paper_A,paper_I}
proposes modifying physical law in localized spacetime regions
by ``muting'' prime channels: localizing
$\mathbb{Z} \to \Zp$. This produces regions where the weak
energy condition is violated, enabling warp-like metrics.

A natural fear: if the modified vacuum has lower energy, could
the bubble expand uncontrollably and replace the normal vacuum
everywhere, destroying the universe? This is the
\emph{vacuum decay catastrophe}, first analyzed by
Coleman and De Luccia~\cite{ColemanDeLuccia1980}.

This paper provides a rigorous answer: \textbf{no}. The
mechanism that prevents universal collapse is fundamentally
different from mere energetic stability --- it is
\emph{topological}.

% ============================================================================
\section{The threat: vacuum decay}
% ============================================================================

In standard QFT, a metastable (``false'') vacuum can decay to
a lower-energy (``true'') vacuum via quantum tunneling.
The decay rate per unit volume is~\cite{ColemanDeLuccia1980}
\begin{equation}
  \Gamma/V = A \exp(-B/\hbar),
  \label{eq:cdl}
\end{equation}
where $B$ is the Euclidean bounce action. In the thin-wall
approximation, $B = 27\pi^2 \sigma^4 / (2\epsilon^3)$,
with $\sigma$ the wall surface tension and $\epsilon$ the
vacuum energy difference.

For arithmetic vacuum engineering: the modified vacuum
$\Spec(\Zp)$ has energy density
$\rho_{\neg p} = \rho_P \zeta_{\neg p}(-3)
= \rho_P (1-p^3)/120 < 0$, which is \emph{lower} than the
normal vacuum $\rho = \rho_P \zeta(-3) = \rho_P/120 > 0$.

In standard CDL theory, a lower-energy vacuum \emph{would}
expand and consume the universe. This is the threat.

% ============================================================================
\section{The topological no-go theorem}
% ============================================================================

\begin{theorem}[Topological stability]
\label{thm:nogo}
If $K_1(\mathbb{Z}) \neq K_1(\Zp)$, no Coleman-De Luccia
bounce solution connecting the $\Spec(\mathbb{Z})$ vacuum to
the $\Spec(\Zp)$ vacuum exists. The nucleation rate is
$\Gamma = 0$ exactly.
\end{theorem}

\begin{proof}
The CDL bounce is a finite-action path
$\gamma: [0,1] \to \mathcal{M}$ in field configuration space
$\mathcal{M}$, with $\gamma(0)$ in the false vacuum sector
and $\gamma(1)$ in the true vacuum sector. The K-theory
invariant $K_1$ is a continuous function on $\mathcal{M}$
(it is defined via homotopy classes of maps, which are
continuous invariants). Along any continuous path $\gamma$,
$K_1(\gamma(t))$ is constant.

Since $K_1(\mathbb{Z}) = \mathbb{Z}/2 \neq \mathbb{Z}/2 \times
\mathbb{Z} = K_1(\Zp)$, the endpoints lie in different
connected components of $\mathcal{M}$. No continuous path
connects them. Therefore no bounce solution exists and
$\Gamma = A \exp(-\infty) = 0$.
\end{proof}

\textbf{Key distinction from standard vacuum decay.}
In conventional QFT, the two vacua are distinguished only by
energy: they have the same topological type. The CDL barrier
is energetic and can be tunneled through. In arithmetic vacuum
engineering, the two vacua have different \emph{topological types}
(different $K_1$). This barrier is structural, not energetic,
and cannot be tunneled through by any quantum process.

\textbf{Analogy.} An energetic barrier is like a mountain:
high, but a tunnel can be dug through it. A topological barrier
is like the distinction between the inside and outside of a
closed surface: no tunnel connects them without puncturing
the surface (changing the topology).

% ============================================================================
\section{Five safety mechanisms}
% ============================================================================

\begin{table}[h]
\centering
\caption{Safety mechanisms for arithmetic vacuum engineering.}
\label{tab:safety}
\small
\begin{tabular}{@{}clll@{}}
\toprule
\# & Mechanism & Principle & Sufficiency \\
\midrule
1 & Topological protection & $K_1$ discrete $\to$ no CDL bounce & \textbf{Sufficient alone} \\
2 & Wall radiation & $T_{\mathrm{wall}} > 0$ $\to$ energy dissipation & Backup only \\
3 & Planck minimum size & $\delta_{\min} \sim \lP \log p$ & Prevents accidents \\
4 & Casimir precedent & 13.8 Gyr of safe vacuum modification & Historical evidence \\
5 & Parity reversibility & $(-1)^{|S|}$ sign law $\to$ Ctrl+Z & Operational safety \\
\bottomrule
\end{tabular}
\end{table}

\textbf{Honest assessment.}
Mechanism~1 (topological protection) is the \emph{only} mechanism
that is individually sufficient to prevent vacuum decay.
Numerical computation shows that for large bubbles, the wall
radiation power (Mechanism~2) is negligible compared to the
volume energy gain: $P_{\mathrm{rad}} / P_{\mathrm{gain}} \sim 10^{-100}$
for meter-scale bubbles. If Mechanism~1 failed (which we argue
it cannot), Mechanism~2 alone would \emph{not} prevent expansion.

This means the safety case rests entirely on
Theorem~\ref{thm:nogo}. The theorem is mathematically rigorous
given the assumption that $K_1$ is a well-defined invariant of the
physical vacuum. This assumption is validated if the vacuum
has $\Spec(\mathbb{Z})$ structure (the premise of the entire
program). If the premise is wrong, the engineering is impossible
anyway, so the safety concern is moot.

% ============================================================================
\section{Counter-arguments}
% ============================================================================

\textbf{Q1: Can finite temperature break topological protection?}

Temperature is a continuous perturbation. $K_1$, as a homotopy
invariant, is stable under continuous perturbations. Protection
fails only above the ``topological gap'' $\Delta \sim E_P$
(Planck energy). No physical environment reaches this
temperature outside the Big Bang singularity.

\textbf{Q2: Can quantum tunneling change topology?}

CDL tunneling is itself formulated as a path integral over
\emph{continuous} field configurations. If no continuous path
connects the two vacua (Theorem~\ref{thm:nogo}), the tunneling
amplitude is zero. Topology change requires a non-perturbative
process beyond the path integral --- such processes may exist
in quantum gravity but are suppressed by
$\exp(-S_{\mathrm{grav}}) \sim \exp(-1/G)$, which is negligibly
small.

\textbf{Q3: Could cosmic rays create bubbles?}

The highest-energy cosmic ray ($\sim 10^{20}$\,eV) is
$\sim 10^8$ times below the Planck energy. Moreover, energy alone
is insufficient: Theorem~\ref{thm:nogo} shows the barrier is
topological, not energetic.

\textbf{Q4: What about black hole interiors?}

Near the singularity, Planck-scale physics may allow topology
change. However, any such change is confined behind the event
horizon and cannot affect the exterior universe.

\textbf{Q5: Can an intentional bubble become uncontrollable?}

The parity law (Mechanism~5) provides an operational kill switch:
muting one additional prime restores the original sign of all
physical corrections. Furthermore, terminating the external
energy supply that maintains the bubble allows wall radiation
(Mechanism~2) to dissipate its energy, though on a potentially
very long timescale for large bubbles.

% ============================================================================
\section{Sandbox: BEC protocol}
% ============================================================================

As a precautionary measure, we propose testing all arithmetic
vacuum modifications in a BEC analog gravity system before
any attempt at real spacetime engineering.

\textbf{Protocol:}
(1)~Create $^{87}$Rb BEC in optical trap.
(2)~Apply periodic optical lattice with period $p$ to create
an ``arithmetic vacuum'' region within the BEC.
(3)~Image the bubble boundary in situ.
(4)~Measure: does the bubble expand, contract, or remain stable?
(5)~Verify parity law: add a second prime modulation and confirm
restoration.
(6)~Scale the bubble size and measure radiation.

If the bubble expands uncontrollably in the BEC, only the
condensate is destroyed --- the universe is unaffected. This
provides a fail-safe testing environment for all subsequent
engineering.

\textbf{Cost:} \$50K (additional optics for existing BEC lab).
\textbf{Timeline:} 6--12 months.

% ============================================================================
\section{Conclusion}
% ============================================================================

The concern that arithmetic vacuum engineering could trigger
universal vacuum collapse is theoretically unfounded. The central
safety guarantee is topological: the $K_1$ invariant difference
between $\Spec(\mathbb{Z})$ and $\Spec(\Zp)$ prevents the
existence of any Coleman-De Luccia bounce solution, making the
nucleation rate exactly zero (Theorem~\ref{thm:nogo}).

We are honest that this is the \emph{only} individually sufficient
mechanism: wall radiation and other backup mechanisms are
inadequate for large bubbles. The safety case therefore depends
on the validity of the topological no-go theorem, which in turn
depends on the premise that the vacuum has $\Spec(\mathbb{Z})$
structure. If this premise is correct (as the evidence in
the companion papers suggests), the engineering is safe.
If the premise is wrong, the engineering is impossible.
Either way, the universe is not at risk.

As a precaution, we recommend testing all operations in a BEC
analog gravity sandbox before deployment.

\begin{thebibliography}{15}
\bibitem{paper_A} Wright Brothers, companion paper A (2026).
\bibitem{paper_I} Wright Brothers, companion paper I (2026).
\bibitem{ColemanDeLuccia1980} S.~Coleman and F.~De Luccia, Phys. Rev. D \textbf{21}, 3305 (1980).
\end{thebibliography}

\end{document}
